\documentclass[12pt]{article}
\usepackage{amsmath, amssymb, amsthm}
\usepackage[margin=1in]{geometry}
\usepackage{algorithm, algorithmicx, algpseudocode}

\newtheorem{theorem}{Theorem}
\newtheorem{lemma}[theorem]{Lemma}

\title{Problem 144: Investigating Multiple Reflections of a Laser Beam}
\author{Project Euler Solution}
\date{}

\begin{document}
\maketitle

\section{Problem Statement}
A laser beam enters an elliptical mirror defined by $4x^2 + y^2 = 100$ and reflects internally. The beam enters at $(0, 10.1)$ and first hits at $(1.4, -9.6)$. The exit is at $-0.01 \le x \le 0.01$ on the top. How many reflections occur before exit?

\section{Mathematical Foundation}

\begin{theorem}[Ellipse normal]
For the ellipse $4x^2 + y^2 = 100$, the outward normal at $(x_0, y_0)$ is proportional to $(4x_0, y_0)$.
\end{theorem}

\begin{proof}
The ellipse is the level set $F(x,y) = 4x^2 + y^2 - 100 = 0$. The gradient $\nabla F = (8x, 2y) \propto (4x, y)$ is perpendicular to the level curve and outward-pointing.
\end{proof}

\begin{theorem}[Reflection formula]
Given incoming direction $\mathbf{d}$ and surface normal $\mathbf{n}$:
\[
\mathbf{d}' = \mathbf{d} - 2\frac{\mathbf{d} \cdot \mathbf{n}}{|\mathbf{n}|^2}\,\mathbf{n}
\]
\end{theorem}

\begin{proof}
Decompose $\mathbf{d} = \mathbf{d}_\parallel + \mathbf{d}_\perp$ where $\mathbf{d}_\parallel = \frac{\mathbf{d} \cdot \mathbf{n}}{|\mathbf{n}|^2}\mathbf{n}$. The law of reflection reverses the normal component while preserving the tangential component: $\mathbf{d}' = \mathbf{d}_\perp - \mathbf{d}_\parallel = \mathbf{d} - 2\mathbf{d}_\parallel$.
\end{proof}

\begin{theorem}[Ray--ellipse intersection]
A ray from $(x_0, y_0)$ on the ellipse in direction $(d_x, d_y)$ hits the ellipse again at parameter:
\[
t = -\frac{2(4x_0 d_x + y_0 d_y)}{4d_x^2 + d_y^2}
\]
\end{theorem}

\begin{proof}
Substituting $(x,y) = (x_0 + td_x, y_0 + td_y)$ into $4x^2 + y^2 = 100$ and using $4x_0^2 + y_0^2 = 100$:
\[
t^2(4d_x^2 + d_y^2) + 2t(4x_0 d_x + y_0 d_y) = 0
\]
The nonzero root gives the stated formula.
\end{proof}

\begin{lemma}
The parameter $t$ is nonzero after each reflection, ensuring the ray always reaches a new point on the ellipse.
\end{lemma}

\begin{proof}
The numerator $2(\mathbf{n} \cdot \mathbf{d})$ vanishes only for tangent directions. After reflection, $\mathbf{d}'$ has nonzero normal component, so $t \neq 0$.
\end{proof}

\section{Algorithm}

\begin{algorithmic}[1]
\State $(x, y) \gets (1.4, -9.6)$; $(d_x, d_y) \gets (1.4, -19.7)$; $\text{count} \gets 0$
\Repeat
    \State $\text{count} \gets \text{count} + 1$
    \State $\mathbf{n} \gets (4x, y)$
    \State $\mathbf{d} \gets \mathbf{d} - 2\frac{\mathbf{d} \cdot \mathbf{n}}{|\mathbf{n}|^2}\,\mathbf{n}$ \Comment{Reflect}
    \State $t \gets -2(4x\,d_x + y\,d_y)/(4d_x^2 + d_y^2)$ \Comment{Next intersection}
    \State $(x, y) \gets (x + t\,d_x,\; y + t\,d_y)$
\Until{$|x| \leq 0.01$ and $y > 0$}
\State \Return count
\end{algorithmic}

\section{Complexity Analysis}

\begin{itemize}
\item \textbf{Time:} $O(R)$ where $R = 354$ is the number of reflections. Each step is $O(1)$.
\item \textbf{Space:} $O(1)$.
\end{itemize}

\section{Answer}

\[
\boxed{354}
\]

\end{document}
